\documentclass[instructions.tex]{subfiles}
\begin{document}
 
\chapter{Du respect dans les églises.}
 
\primo Ma maison est une maison de prière ; n’en faites pas une retraite de voleurs, disait Jésus-Christ aux juifs. Nos églises sont bien plus saintes, que le temple des juifs; elles sont cependant plus indignement profanées. On n’en fait pas seulement une retraite de voleurs, mais une retraite d’impies, de voluptueux, de sacriléges, par les discours qu’on y tient, par les postures qu’on y prend, par les vanités qu’on y affecte, par les regards qu’on y jette, par les pensées et les désirs dont on se souille, par les rendez-vous qu’on y donne, par les sacremens qu’on y profane.
 
On entre avec respect dans la maison des grands, qui sont peut-être de grands pécheurs, et on n’a point de honte d’entrer sans respect dans la maison de Dieu. Les anges y tremblent : J’ai vu, dit saint Chrysostome, ces esprits célestes, pendant les saints mystères, se prosterner devant Jésus-Christ et l’adorer; pendant qu’on voit des libertins, de jeunes étourdis, des filles effrontées, rire, sourire, se dissiper comme sur une place publique.
 
En présence des rois de la terre, on garde le silence, on n’ose y faire paraître la moindre légèreté; et l’on voit dans les églises, en présence du Roi du ciel, des gens pleins de fierté, fléchir à peine les genoux pendant le sacrifice, n’y assister que pour se faire remarquer, pour mépriser le peuple et insulter à Dieu!
 
Où allez vous, dit un jour saint Ambroise à une dame superbement parée?... Je vais à l’église, répondit-elle... On dirait bien plutôt, répliqua le saint pasteur, que vous allez à la comédie ou à la danse; allez, retirez-vous, femme vaine et pécheresse ; venez dans le lieu saint pour y pleurer vos crimes, non pas pour insulter à Jésus-Christ et scandaliser.
 
On n’oserait mener un chien dans le cabinet d’un seigneur, et l’on voit dans la maison de Dieu ces vils animaux distraire la piété des assistans, troubler les offices divins, sans que l’on ose rien dire. On estime un domestique qui empêche qu’on ne fasse des indécences dans la maison de son maître, et l’on blâme souvent un ministre de Dieu, un pasteur qui s’oppose aux irrévérences qu’on commet en présence de Jésus-Christ.
 
Un père se croit offensé, si l’on insulte à son enfant; et ce père voit son enfant à l’église insulter à Jésus-Christ par ses badineries, et ne dit rien. Un maître veut que ses écoliers le respectent; et il voit de petits impies sans respect devant Jésus-Christ, et ne les châtie point. Un grand veut être honoré de ses subordonnés; et lui-même, par ses airs fastueux, déshonore Jésus-Christ en leur présence, et souffre, sans y mettre ordre, qu’ils le déshonorent aussi par leurs irrévérences.
 
\secundo Quoi! les païens tremblent devant leurs faux dieux, et les chrétiens sont sans respect devant le Tout-Puissant! Est-il possible qu’on ait moins de retenue dans la maison du vrai Dieu, que les païens n’en ont dans le temple d’une idole; et qu’on ait moins de respect pour le Très-Haut, que les infidèles n’en ont pour des statues? Dans quel endroit Dieu sera-t-il à couvert de nos insultes, s’il est outragé dans sa propre maison?
 
O mon Dieu! de telles profanations seront-elles toujours impunies? Je m’étonne, dit saint Chrysostome, que la foudre n’écrase pas ces téméraires. Quel compte n’en rendront pas à Dieu ceux qui ont l’autorité, et qui n’empêchent pas de tels abus! Si nous n’avons pas assez de pouvoir pour les empêcher, ayons du moins assez de zèle et de foi pour les déplorer.
 
Tâchons de dédommager notre Sauveur de ces outrages; gémissons sur l’aveuglement de tant de chrétiens, qui lui font injure dans le temps même qu’il se sacrifie pour eux. Imitons la sainte Vierge et le disciple du Sauveur, qui lui compatissaient pendant qu’on le crucifiait, et lui tenaient compagnie pendant qu’il était abandonné de tous.
 
\section{Résolutions}
 
\primo Quand j’entrerai dans une église, je me rappellerai qu’elle est la maison de Dieu.
 
\secundo Je n’y parlerai que dans la nécessité, et brièvement.
 
\tertio J’y tiendrai toujours une contenance respectueuse, et je m’y occuperai à prier le Seigneur, à entendre sa parole sainte, et à traiter avec lui de la grande affaire de mon salut, sans laisser égarer mes yeux ni mes pensées sur des objets dangereux ou seulement inutiles. 
\prayer{Mon Dieu, pardonnez-moi les péchés que j’ai commis dans vos saints temples et jusqu’au pied de vos autels.}
 
\end{document}
